% Monte Carlo: throw random points at the unit square and count how many land
% inside the quarter circle. The ratio approaches pi/4, and every trial is
% independent -- which is what makes it embarrassingly parallel.
\documentclass[dvisvgm,border=3pt]{standalone}
\input{preamble.tex}
\begin{document}
\begin{tikzpicture}[x=1cm,y=1cm]

  \def\R{4.2}
  \draw[draw=blu,line width=1.1pt] (0,0) rectangle (\R,\R);
  % quarter circle centred on the origin: start the arc at (R,0) so its centre
  % lands at (0,0) rather than outside the square
  \draw[draw=uibkorange,line width=1.1pt]
    (\R,0) arc[start angle=0,end angle=90,radius=\R];

  % a fixed sample so the picture is reproducible; inside the arc = orange
  \foreach \x/\y in {0.42/0.55, 1.10/0.38, 0.78/1.42, 1.85/0.62, 2.40/1.05,
                     1.35/2.15, 0.55/2.70, 2.95/1.55, 0.95/3.30, 2.05/2.55,
                     3.35/0.85, 1.62/1.28, 2.62/2.05, 0.30/1.95, 3.05/2.35,
                     1.18/0.92, 2.25/3.15, 0.68/3.85, 3.62/1.42, 1.95/1.78,
                     2.78/0.45, 0.22/3.15, 1.48/3.52, 3.18/1.95, 2.48/0.82}{
    \fill[uibkorange] (\x,\y) circle (0.075);}
  \foreach \x/\y in {3.95/2.30, 3.55/3.05, 2.95/3.45, 3.80/3.70, 2.55/3.85,
                     3.35/3.35, 4.02/1.60, 3.72/2.72}{
    \fill[blu] (\x,\y) circle (0.075);}

\end{tikzpicture}
\end{document}
